% =============================================================================
%  lec20.tex —— 第 20 讲 拉直与反拉直（底本 PDF 第 153–160 页）
% =============================================================================
\lecture{20}{拉直与反拉直}
\begin{center}\small\color{black!55}（底本 PDF 第 153–160 页）\end{center}

本讲对（小）$\infty$-范畴 $\mathrm{D}$，说明下述两者之间的一个典范等价：
\begin{enumerate}[label=(\roman*),leftmargin=2.4em,itemsep=0.1em]
\item $\mathrm{D}$ 上笛卡尔纤维化构成的 $\infty$-范畴；
\item 函子 $\mathrm{D}^{\op}\to\Cat_\infty$ 构成的 $\infty$-范畴。
\end{enumerate}
构造 (i)$\Rightarrow$(ii) 称为\textbf{拉直}，而 (ii)$\Rightarrow$(i) 称为\textbf{反拉直}。对普通范畴，相应的等价由 Grothendieck 证明，并以\textbf{Grothendieck 构造}之名著称。对 $\infty$-范畴，人们称之为 \textbf{Grothendieck–Lurie 构造}。

\yizhu{straightening 与 unstraightening 分别译为“拉直”与“反拉直”。旧译“平展/去平展”今已不取，因为“平展”通行用作 étale 的译名，两者含义不同。\textbf{[来源：B]}}

\section{传输函子}

\nitem\label{L20-1.1} 本节说明如何构造笛卡尔纤维化各纤维之间的函子。

设 $p\colon \mathrm{C}\to[1]$ 是 $\infty$-范畴之间的笛卡尔纤维化。考虑 $p$ 的\textbf{截面}构成的 $\infty$-范畴\footnote{下面的纤维积没有歧义，因为 $\Fun([1],\mathrm{C})\to\Fun([1],[1])$ 是范畴纤维化。}：
\[
\operatorname{Sect}_{/[1]}(\mathrm{C})\defeq \Fun([1],\mathrm{C})\times_{\Fun([1],[1])}\{\id\}.
\]
设
\[
\operatorname{Sect}^{\mathrm{Cart}}_{/[1]}(\mathrm{C})\subseteq\operatorname{Sect}_{/[1]}(\mathrm{C})
\]
是满子 $\infty$-范畴，其对象由 $\mathrm{C}$ 中的 $p$-笛卡尔箭头给出。

\begin{proposition}\label{L20-1.2}
设 $p\colon \mathrm{C}\to[1]$ 是 $\infty$-范畴之间的笛卡尔纤维化。则函子
\[
\mathrm{ev}_1\colon \operatorname{Sect}^{\mathrm{Cart}}_{/[1]}(\mathrm{C})\to\mathrm{C}_1
\]
是 $\infty$-范畴之间的等价。
\end{proposition}

\begin{proof}[证明梗概]
该函子本质满，因为 $p$-笛卡尔箭头足够多。剩下的只需证明它全忠实。对 $\operatorname{Sect}^{\mathrm{Cart}}_{/[1]}(\mathrm{C})$ 中的对象 $x\xrightarrow{f}y$ 与 $x'\xrightarrow{f'}y'$，有
\[
\mathrm{Maps}(f,f')\simeq \mathrm{Maps}_{\mathrm{C}_0}(x,x')\times_{\mathrm{Maps}_{\mathrm{C}}(x,y')}\mathrm{Maps}_{\mathrm{C}_1}(y,y').
\]
由于 $f'$ 是局部 $p$-笛卡尔的，函子 $\mathrm{Maps}_{\mathrm{C}_0}(x,x')\to\mathrm{Maps}_{\mathrm{C}}(x,y')$ 是等价。于是函子 $\mathrm{Maps}(f,f')\to\mathrm{Maps}_{\mathrm{C}_1}(y,y')$ 是等价，此即所求。
\end{proof}

\begin{construction}\label{L20-1.3}
设 $p\colon \mathrm{C}\to\mathrm{D}$ 是 $\infty$-范畴之间的笛卡尔纤维化。由命题~\ref{L20-1.2}，对 $\mathrm{D}$ 中任意态射 $f\colon x\to y$，有等价
\[
\operatorname{Sect}^{\mathrm{Cart}}_{[1]}(\mathrm{C}\times_{\mathrm{D}}[1])\xrightarrow{\ \simeq\ }\mathrm{C}_y,
\]
其中 $[1]\to\mathrm{D}$ 是由 $f$ 给出的函子。选取上述等价的一个逆，我们得到函子
\[
\mathrm{C}_y\xrightarrow{\ \simeq\ }\operatorname{Sect}^{\mathrm{Cart}}_{[1]}(\mathrm{C}\times_{\mathrm{D}}[1])\xrightarrow{\mathrm{ev}_0}\mathrm{C}_x.
\]
我们把这个函子记为 $f^{\dagger}\colon\mathrm{C}_y\to\mathrm{C}_x$，并称之为沿 $f$ 的\textbf{反变传输函子}。

对偶地，对余笛卡尔纤维化 $p\colon \mathrm{C}\to\mathrm{D}$，我们定义\textbf{协变传输函子} $f_{\dagger}\colon\mathrm{C}_x\to\mathrm{C}_y$。
\end{construction}

\yizhu{原文用剑号 $f^{\dagger}$、$f_{\dagger}$ 记传输函子，而 Lurie 的 HTT 中对应记号是 $f^{*}$、$f_{*}$（分别对应拉回与前推）。此处保留原记号，以免与拉回函子混淆。\textbf{[来源：B]}}

\nitem\label{L20-1.4} 由构造，$f^{\dagger}$ 在一个可缩的选择空间意义下良定义。

\begin{exercise}\label{L20-1.5}
设 $p\colon \mathrm{C}\to\mathrm{D}$ 是 $\infty$-范畴之间的笛卡尔纤维化，$f,f'\colon x\to y$ 是 $\mathrm{D}$ 中彼此同伦的态射。证明
\[
f^{\dagger}\simeq(f')^{\dagger}.
\]
\end{exercise}

\begin{exercise}\label{L20-1.6}
设 $p\colon \mathrm{C}\to\mathrm{D}$ 是 $\infty$-范畴之间的笛卡尔纤维化，$x\xrightarrow{f}y\xrightarrow{g}z$ 是 $\mathrm{D}$ 中的一条链。证明
\[
f^{\dagger}\circ g^{\dagger}\simeq(g\circ f)^{\dagger}.
\]
\end{exercise}

\begin{exercise}\label{L20-1.7}
注意构造 $f\mapsto f^{\dagger}$ 对拟范畴之间的任意局部笛卡尔纤维化 $p\colon\mathrm{C}\to\mathrm{D}$ 都有意义。证明在此情形下，存在典范自然变换
\[
f^{\dagger}\circ g^{\dagger}\to(g\circ f)^{\dagger}.
\]
\end{exercise}

\begin{exercise}\label{L20-1.8}
设 $p\colon \mathrm{C}\to\mathrm{D}$ 既笛卡尔又余笛卡尔的纤维化。证明对 $\mathrm{D}$ 中任意态射 $f\colon x\to y$，$f_{\dagger}$ 是 $f^{\dagger}$ 的左伴随。
\end{exercise}

\nitem\label{L20-1.9} 由上述练习，对 $\infty$-范畴之间的任意笛卡尔纤维化 $p\colon \mathrm{C}\to\mathrm{D}$，我们有典范函子
\[
\mathrm{D}^{\op}\to\mathrm{h}\Cat_\infty,\qquad x\mapsto\mathrm{C}_x,\qquad (x\xrightarrow{f}y)\mapsto f^{\dagger}.
\]
类似地，对任意余笛卡尔纤维化 $p\colon \mathrm{C}\to\mathrm{D}$，我们有典范函子 $\mathrm{D}\to\mathrm{h}\Cat_\infty$。

\section{带标记单纯集}

\nitem\label{L20-2.1} 设 $p\colon \mathrm{C}\to\mathrm{D}$ 是 $\infty$-范畴之间的笛卡尔纤维化，我们希望把上述函子 $\mathrm{D}^{\op}\to\mathrm{h}\Cat_\infty$ 提升为函子 $\mathrm{D}^{\op}\to\Cat_\infty$。此外，我们还希望这样的构造对 $p$ 具有函子性。为了确切描述这个目标，我们引入下述 $\infty$-范畴。

\begin{construction}\label{L20-2.2}
对 $\mathrm{D}\in\Cat_\infty$，设
\[
\mathrm{Cart}_{/\mathrm{D}}\subseteq(\Cat_\infty)_{/\mathrm{D}}
\]
是这样的子 $\infty$-范畴：
\begin{enumerate}[label=(\roman*),leftmargin=2.4em,itemsep=0.1em]
\item 对象是 $\mathrm{D}$ 上的笛卡尔纤维化；
\item 态射是定义在 $\mathrm{D}$ 上、保持笛卡尔箭头的函子。
\end{enumerate}
\end{construction}

\begin{remark}\label{L20-2.3}
更确切地说，我们把 $(\Cat_\infty)_{/\mathrm{D}}$ 实现为拟范畴，并把 $\mathrm{Cart}_{/\mathrm{D}}$ 定义为具有指定 $1$-骨架的极大单纯子集。
\end{remark}

\begin{warning}\label{L20-2.4}
嵌入 $\mathrm{Cart}_{/\mathrm{D}}\subseteq(\Cat_\infty)_{/\mathrm{D}}$ 不是全忠实的。它是 $1$-全忠实的，这意味着它在映射 $\infty$-群胚之间诱导全忠实函子。
\end{warning}

\nitem\label{L20-2.5} 我们的目标是构造一个等价
\begin{equation}\label{L20-eq2.1}
\mathrm{Cart}_{/\mathrm{D}}\simeq\Fun(\mathrm{D}^{\op},\Cat_\infty).
\end{equation}
回忆：已知的 $\Cat_\infty$ 的构造只有两种，都以模型范畴为基础。即：$\Cat_\infty$ 等价于模型范畴 $\Set^{\mathrm{Joyal}}_\Delta$ 底层的 $\infty$-范畴，而后者可用下述等价的方式构造：
\begin{itemize}[leftmargin=2.2em,itemsep=0.1em]
\item 以 $\infty$-范畴的方式把 $\Set_\Delta$ 中的范畴等价反演。
\item 把满子范畴 $\mathcal{Q}\mathrm{Cat}\simeq(\Set^{\mathrm{Joyal}}_\Delta)^\circ$ 以单纯集充实，并取单纯脉函子 $\mathfrak{N}_\bullet$。
\end{itemize}
因此，构造等价 \eqref{L20-eq2.1} 的唯一合理策略，是把它实现为模型范畴之间的 Quillen 等价，而这两个模型范畴底层的 $\infty$-范畴分别是 $\mathrm{Cart}_{/\mathrm{D}}$ 与 $\Fun(\mathrm{D}^{\op},\Cat_\infty)$。

\nitem\label{L20-2.6} 我们先为 $\mathrm{Cart}_{/\mathrm{D}}$ 寻找一个模型。把 $\mathrm{D}$ 实现为拟范畴 $\mathcal{D}$ 后，人们可能想在 $(\Set_\Delta)_{/\mathcal{D}}$ 上找一个能达到上述目标的模型结构。然而这个希望并不合理：即便它的双纤维对象——即 $\mathcal{D}$ 上的笛卡尔纤维化——是正确的，这些对象之间的态射也将是定义在 $\mathcal{D}$ 上的所有函子，而不是那些保持笛卡尔箭头的函子。这说明我们需要修改 $(\Set_\Delta)_{/\mathcal{D}}$，使得态射自动保持某些特殊的箭头。这引出下述定义。

\begin{definition}\label{L20-2.7}
设 $\Set^+_\Delta$ 是如下定义的普通范畴。
\begin{itemize}[leftmargin=2.2em,itemsep=0.1em]
\item 对象是二元组 $(X,E)$，其中 $X$ 是单纯集，$E\subseteq X_1$ 是包含全部退化边的边集。
\item 从 $(X,E)$ 到 $(X',E')$ 的态射是单纯集之间的态射 $X\to X'$，它把 $E$ 送入 $E'$。
\end{itemize}
我们称上述二元组 $(X,E)$ 为\textbf{带标记单纯集}，并称 $E$ 中的边为\textbf{带标记的边}。
\end{definition}

\begin{example}\label{L20-2.8}
对任意单纯集 $S$，
\begin{itemize}[leftmargin=2.2em,itemsep=0.1em]
\item 设 $S^{\sharp}\in\Set^+_\Delta$ 是这样的对象：$S$ 中所有边都带标记；
\item 设 $S^{\flat}\in\Set^+_\Delta$ 是这样的对象：$S$ 中只有退化边带标记。
\end{itemize}
\end{example}

\begin{exercise}\label{L20-2.9}
证明函子 $\Set_\Delta\to\Set^+_\Delta$、$S\mapsto S^{\sharp}$ 是遗忘函子 $\oblv\colon\Set^+_\Delta\to\Set_\Delta$ 的右伴随。对偶地，$S\mapsto S^{\flat}$ 是遗忘函子的左伴随。
\end{exercise}

\begin{definition}\label{L20-2.10}
设 $S$ 是单纯集。定义
\[
(\Set^+_\Delta)_{/S}\defeq(\Set^+_\Delta)_{/S^{\sharp}},
\]
并称之为 $S$ 上的\textbf{带标记单纯集范畴}。
\end{definition}

\begin{example}\label{L20-2.11}
设 $X\to S$ 是单纯集之间的笛卡尔纤维化，$E\subseteq X_1$ 是全部笛卡尔边构成的子集。则 $X^{\natural}\defeq(X,E)$ 是 $(\Set^+_\Delta)_{/S}$ 中的对象。
\end{example}

\nitem\label{L20-2.12} 我们将要在 $(\Set^+_\Delta)_{/S}$ 上构造一个模型结构，使得
\begin{itemize}[leftmargin=2.2em,itemsep=0.1em]
\item 上纤维化是单态射；
\item 纤维对象是 $X^{\natural}$，其中 $X\to S$ 是笛卡尔纤维化。
\end{itemize}
正如第 5 讲第 4 节所解释的，只要我们知道如何计算所需 $\infty$-范畴 $(\Set^+_\Delta)_{/S}[W^{-1}]$ 中的映射空间，这些要求就确定了该模型结构。

\begin{construction}\label{L20-2.13}
对 $Y,Z\in\Set^+_\Delta$，考虑由下述泛性质定义的单纯集 $\Hom^{\sharp}(Y,Z)$：对任意 $K\in\Set_\Delta$，
\[
\Hom_{\Set_\Delta}(K,\Hom^{\sharp}(Y,Z))\simeq\Hom_{\Set^+_\Delta}(K^{\sharp}\times Y,Z).
\]
对 $Y,Z\in(\Set^+_\Delta)_{/S}$，设
\[
\Hom^{\sharp}_S(Y,Z)\subseteq\Hom^{\sharp}(Y,Z)
\]
是满单纯子集，其顶点由定义在 $S$ 上的态射 $Y\to Z$ 给出。
\end{construction}

\begin{exercise}\label{L20-2.14}
设 $Y\in(\Set^+_\Delta)_{/S}$，$Z\to S$ 是笛卡尔纤维化。证明 $\Hom^{\sharp}_S(Y,Z^{\natural})$ 是 Kan 复形。请说服自己：这个 Kan 复形建模的是定义在 $S$ 上、把 $Y$ 中带标记的边送入 $Z$ 中笛卡尔箭头的函子 $Y\to Z$ 构成的空间。
\end{exercise}

\begin{exercise}\label{L20-2.15}
若改用 $\Hom^{\flat}(Y,Z^{\natural})$，会发生什么？
\end{exercise}

\begin{definition}\label{L20-2.16}
$(\Set^+_\Delta)_{/S}$ 中的态射 $Y\to Y'$ 称为\textbf{笛卡尔等价}，若对任意笛卡尔纤维化 $Z\to S$，函子
\[
\Hom^{\sharp}_S(Y',Z^{\natural})\to\Hom^{\sharp}_S(Y,Z^{\natural})
\]
是 Kan 复形之间的弱等价。
\end{definition}

\begin{remark}\label{L20-2.17}
由 HTT.3.1.3.3，若改用 $\Hom^{\flat}(Y,Z^{\natural})$，我们得到相同的概念。
\end{remark}

\begin{theorem}\label{L20-2.18}（HTT.3.1.3.7, 3.1.4.1, 3.1.4.4）
设 $S$ 是单纯集。在 $(\Set^+_\Delta)_{/S}$ 上存在一个左真组合模型结构，称为\textbf{笛卡尔模型结构}，使得
\begin{itemize}[leftmargin=2.2em,itemsep=0.1em]
\item[(C)] 上纤维化是单态射；
\item[(W)] 弱等价是笛卡尔等价。
\end{itemize}
此外，
\begin{enumerate}[label=(\arabic*),leftmargin=2.4em,itemsep=0.1em]
\item 一个对象是纤维的，当且仅当它形如 $X^{\natural}$，其中 $X\to S$ 是笛卡尔纤维化。
\item 该模型结构与由 $\Hom^{\sharp}(-,-)$ 给出的 $\Set^{\mathrm{KQ}}_\Delta$-充实相容\footnote{关于这句话的含义，见第 8 讲第 2 节。}。
\end{enumerate}
\end{theorem}

\nitem\label{L20-2.19} 我们把上述模型范畴记为 $(\Set^+_\Delta)^{\mathrm{Cart}}_{/S}$，并把 $\Set^{\mathrm{KQ}}_\Delta$-充实的模型范畴记为 $(\widetilde{\Set}^+_\Delta)^{\mathrm{Cart}}_{/S}$。对偶地，可以在 $(\Set^+_\Delta)_{/S}$ 上定义余笛卡尔模型结构。

\yizhu{原文以 $\widetilde{\Set}^+_\Delta$ 记以 $\Set^{\mathrm{KQ}}_\Delta$ 充实的带标记单纯集范畴，以区别于普通范畴 $\Set^+_\Delta$；这个波浪号在排印中很小，容易看漏。\textbf{[来源：I]}}

\nitem\label{L20-2.20} 当 $S=\Delta^0$ 时，笛卡尔模型结构与余笛卡尔模型结构重合，因此我们就把所得的模型范畴简记为 $\Set^+_\Delta$。正如人们所期待的，$\Set^+_\Delta$ 的同伦论也等价于 $\Set^{\mathrm{Joyal}}_\Delta$。

\begin{theorem}\label{L20-2.21}（HTT.3.1.5.1）
当 $S=\Delta^0$ 时，伴随
\[
(-)^{\flat}\colon\Set^{\mathrm{Joyal}}_\Delta\rightleftarrows\Set^+_\Delta:\oblv
\]
是 Quillen 等价。
\end{theorem}

\begin{exercise}\label{L20-2.22}
设 $(X,E)\in\Set^+_\Delta$ 是带标记单纯集。证明它在 $\Set^+_\Delta[W^{-1}]\simeq\Cat_\infty$ 中的像等价于 $X[E^{-1}]$。
\end{exercise}

\begin{exercise}\label{L20-2.23}
设 $X\to\Delta^1$ 是拟范畴之间的笛卡尔纤维化。证明
\[
(\widetilde{\Set}^+_\Delta)_{/\Delta^1}\to\widetilde{\Set}^+_\Delta\to\Set^+_\Delta[W^{-1}]\simeq\Cat_\infty
\]
把 $X^{\natural}$ 送到一个等价于 $X_0$ 的 $\infty$-范畴。此外，$X^{\natural}_1\to X^{\natural}$ 的像建模反变传输函子 $X_1\to X_0$。这里 $X^{\natural}_1$ 把所有同构都标为带标记的。
\end{exercise}

\begin{exercise}\label{L20-2.24}
设 $\mathrm{D}$ 是实现在拟范畴 $\mathcal{D}$ 上的 $\infty$-范畴。证明 $(\Set^+_\Delta)^{\mathrm{Cart}}_{/\mathcal{D}}$ 底层的 $\infty$-范畴等价于 $\mathrm{Cart}_{/\mathrm{D}}$。
\end{exercise}

\begin{exercise}\label{L20-2.25}
设 $\mathrm{D}$ 是实现在单纯范畴 $\mathbb{D}$ 上的 $\infty$-范畴。考虑单纯函子 $\mathbb{D}\to\Set^+_\Delta$ 的投影模型范畴\footnote{关于这句话的含义，见第 8 讲第 2 节。}：
\[
\Fun(\mathbb{D},\Set^+_\Delta)_{\mathrm{proj}}.
\]
证明它底层的 $\infty$-范畴等价于 $\Fun(\mathrm{D},\Cat_\infty)$。
\end{exercise}

\section{拉直与反拉直}

\nitem\label{L20-3.1} 设 $\infty$-范畴 $\mathrm{D}$ 同时实现为拟范畴 $\mathcal{D}$ 与单纯范畴 $\mathbb{D}$。由上述练习，为了构造所需的等价，我们只需构造一个 Quillen 等价
\[
(\Set^+_\Delta)^{\mathrm{Cart}}_{/\mathcal{D}}\rightleftarrows\Fun(\mathbb{D}^{\op},\Set^+_\Delta)_{\mathrm{proj}}.
\]
按定义，单纯范畴之间存在弱等价 $\phi\colon\mathfrak{C}(\mathcal{D})\to\mathbb{D}$。

\begin{theorem}\label{L20-3.2}（HTT.3.2.0.1）
设 $\mathcal{D}$ 是单纯集，$\mathbb{D}$ 是单纯范畴。对任意函子 $\phi\colon\mathfrak{C}(\mathcal{D})\to\mathbb{D}$，存在典范 Quillen 伴随
\[
\mathrm{St}^+_{\phi}\colon(\Set^+_\Delta)^{\mathrm{Cart}}_{/\mathcal{D}}\rightleftarrows\Fun(\mathbb{D}^{\op},\Set^+_\Delta)_{\mathrm{proj}}:\mathrm{Un}^+_{\phi},
\]
若 $\phi$ 是弱等价，则它是 Quillen 等价。
\end{theorem}

\begin{corollary}\label{L20-3.3}
设 $\mathrm{D}$ 是任意 $\infty$-范畴。存在典范等价
\[
\mathrm{Cart}_{/\mathrm{D}}\simeq\Fun(\mathrm{D}^{\op},\Cat_\infty).
\]
\end{corollary}

\nitem\label{L20-3.4} 我们来说明 $\mathrm{St}^+_{\phi}$ 的构造要点。先处理 $\mathbb{D}=\mathfrak{C}(\mathcal{D})$ 的情形。

对 $\mathcal{D}$ 上的带标记单纯集 $X\in(\Set^+_\Delta)_{/\mathcal{D}}$，我们希望构造单纯函子
\[
F\colon\mathfrak{C}(\mathcal{D})^{\op}\to\Set^+_\Delta,
\]
使得对 $\mathcal{D}$ 中任意对象 $d$（视为 $\mathfrak{C}(\mathcal{D})$ 中的对象），$F(d)$ 都弱等价于 $X_d$。朴素的构造 $d\mapsto X_d$ 行不通，因为无法把函子性做出来。练习~\ref{L20-2.23} 提示，
\begin{center}
\emph{$F(d)$ 应当把 $X$ 在所有可能的箭头 $d\to d'$ 上的信息编码进来。}
\end{center}
例如，我们希望 $X$ 中位于 $d\to d'$ 上方的笛卡尔箭头，都是 $F(d)$ 中带标记的边。

假设我们找到了实现上述目标的办法，就会得到态射 $F(d')\to F(d)$，它建模沿 $d\to d'$ 的反变传输函子。然而，要得到同伦连贯的复合律，我们意识到 $F(d)$ 应当把 $X$ 在所有可能的链 $d\to d'\to d''\to\cdots$ 上的信息都编码进来。

换言之，我们希望 $F(d)$ 的顶点是链 $x_0\to x_1\to\cdots x_n$，其中 $x_0\in X_d$，并且链的加细可以编码为 $F(d)$ 中的边。下述构造提供了实现这一想法的途径。

\begin{construction}\label{L20-3.5}
设 $p\colon X\to\mathcal{D}$ 是 $\mathcal{D}$ 上的任意单纯集。定义 $\mathbb{D}_X$ 为 $\Cat_\Delta$ 中的下述推出：
\begin{center}
\begin{tikzcd}[column sep=4.6em,row sep=2.8em]
\mathfrak{C}(X) \arrow[r] \arrow[d] & \mathfrak{C}(X^{\triangleright}) \arrow[dd, dotted] \\
\mathfrak{C}(\mathcal{D}) \arrow[d, "\phi"'] & \\
\mathbb{D} \arrow[r, dotted] & \mathbb{D}_X .
\end{tikzcd}
\end{center}
考虑函子
\[
\mathrm{St}_{\phi}X\colon\mathbb{D}^{\op}\to\Set_\Delta,\qquad d\mapsto\Hom_{\mathbb{D}_X}(d,\ast).
\]
我们称它为 $X$ 关于 $\phi$ 的\textbf{拉直}。
\end{construction}

\begin{construction}\label{L20-3.6}
设 $p\colon X\to\mathcal{D}$ 是 $\mathcal{D}$ 上的带标记单纯集。对 $d\to d'$ 上方的一条带标记的边 $x\to x'$，考虑 $2$-单形 $\Delta^2\to X^{\triangleright}$，$0\mapsto x$，$1\mapsto x'$，$2\mapsto\ast$。函子
\[
\mathfrak{C}(\Delta^2)\to\mathfrak{C}(X^{\triangleright})\to\mathfrak{C}(\mathbb{D}_X)
\]
诱导单纯集之间的态射
\[
\Delta^1\simeq\Hom_{\mathfrak{C}(\Delta^2)}(0,2)\to\Hom_{\mathfrak{C}(\mathbb{D}_X)}(d,\ast)\mathrel{=:}\mathrm{St}_{\phi}X(d).
\]
我们把 $\mathrm{St}_{\phi}X(d)$ 中按上述方式得到的边宣告为\textbf{特殊边}，并定义带标记单纯集
\[
\mathrm{St}^+_{\phi}X(d)\defeq(\mathrm{St}_{\phi}X(d),E_d),
\]
使得 $\{E_d\}$ 是包含全部特殊边、且函子地依赖于 $d$ 的最小的一族集合。所得的函子
\[
\mathrm{St}^+_{\phi}X\colon\mathbb{D}^{\op}\to\Set^+_\Delta
\]
与单纯充实相容，并且对 $X$ 有函子性。这就给出了所需的函子
\[
\mathrm{St}^+_{\phi}\colon(\Set^+_\Delta)^{\mathrm{Cart}}_{/\mathcal{D}}\to\Fun(\mathbb{D}^{\op},\Set^+_\Delta)_{\mathrm{proj}}.
\]
\end{construction}

\section{笛卡尔纤维化的基变换}

\begin{proposition}\label{L20-4.1}（HTT.3.2.1.4）
对 $i=0,1$，设 $\mathcal{D}_i$ 是单纯集，$\mathbb{D}_i$ 是单纯范畴。假设在 $\Cat_\Delta$ 中有下述交换图：
\begin{center}
\begin{tikzcd}[column sep=5.6em,row sep=2.8em]
\mathfrak{C}(\mathcal{D}_1) \arrow[r, "\phi_1"] \arrow[d, "\mathfrak{C}(F)"] & \mathbb{D}_1 \arrow[d, "F"] \\
\mathfrak{C}(\mathcal{D}_2) \arrow[r, "\phi_2"] & \mathbb{D}_2 .
\end{tikzcd}
\end{center}
则左 Quillen 函子之间有下述交换图
\begin{center}
\begin{tikzcd}[column sep=8em,row sep=3.6em]
(\Set^+_\Delta)^{\mathrm{Cart}}_{/\mathcal{D}_1} \arrow[r, "\mathrm{St}^+_{\phi_1}"] \arrow[d, "F\circ-"] & \Fun(\mathbb{D}_1^{\op},\widetilde{\Set}^+_\Delta)_{\mathrm{proj}} \arrow[d, "\mathrm{LKE}_{F^{\op}}"] \\
(\Set^+_\Delta)^{\mathrm{Cart}}_{/\mathcal{D}_2} \arrow[r, "\mathrm{St}^+_{\phi_2}"] & \Fun(\mathbb{D}_2^{\op},\widetilde{\Set}^+_\Delta)_{\mathrm{proj}} ,
\end{tikzcd}
\end{center}
而对它们的右伴随也类似：
\begin{center}
\begin{tikzcd}[column sep=8em,row sep=3.6em]
(\Set^+_\Delta)^{\mathrm{Cart}}_{/\mathcal{D}_1} & \Fun(\mathbb{D}_1^{\op},\widetilde{\Set}^+_\Delta)_{\mathrm{proj}} \arrow[l, "\mathrm{Un}^+_{\phi_1}"'] \\
(\Set^+_\Delta)^{\mathrm{Cart}}_{/\mathcal{D}_2} \arrow[u, "-\times_{\mathcal{D}_2}\mathcal{D}_1"'] & \Fun(\mathbb{D}_2^{\op},\widetilde{\Set}^+_\Delta)_{\mathrm{proj}} \arrow[l, "\mathrm{Un}^+_{\phi_2}"'] \arrow[u, "-\circ F^{\op}"'] .
\end{tikzcd}
\end{center}
\end{proposition}

\begin{corollary}\label{L20-4.2}
设 $F\colon\mathrm{D}_1\to\mathrm{D}_2$ 是 $\infty$-范畴之间的函子。存在典范交换图
\begin{center}
\begin{tikzcd}[column sep=7em,row sep=3.4em]
\mathrm{Cart}_{/\mathrm{D}_1} & \Fun(\mathrm{D}_1^{\op},\Cat_\infty) \arrow[l, "\simeq"'] \\
\mathrm{Cart}_{/\mathrm{D}_2} \arrow[u, "-\times_{\mathrm{D}_2}\mathrm{D}_1"'] & \Fun(\mathrm{D}_2^{\op},\Cat_\infty) \arrow[l, "\simeq"'] \arrow[u, "-\circ F^{\op}"'] .
\end{tikzcd}
\end{center}
\end{corollary}

\section{应用}

\nitem\label{L20-5.1} 配备 Grothendieck–Lurie 构造的笛卡尔纤维化理论极其强大。下面是一些例子。

\begin{example}\label{L20-5.2}
对任意 $\infty$-范畴 $\mathrm{D}$，笛卡尔纤维化
\[
\mathrm{ev}_0\colon\Fun(\Delta^1,\mathrm{D})\to\mathrm{D}
\]
对应一个函子 $\mathrm{D}^{\op}\to\Cat_\infty$，它在每个 $d\in\mathrm{D}$ 处的值等价于 $\mathrm{D}_{d/}$。这为（余）切片构造提供了函子性。
\end{example}

\begin{example}\label{L20-5.3}
对 $\infty$-范畴之间的任意函子 $F\colon\mathrm{C}\to\mathrm{D}$。可以证明
\[
\mathrm{ev}_0\colon\Fun(\Delta^1,\mathrm{D})\times_{\Fun(\{1\},\mathrm{D})}\mathrm{C}\to\mathrm{D}
\]
是笛卡尔纤维化，它给出一个函子 $\mathrm{D}^{\op}\to\Cat_\infty$，其在 $d\in\mathrm{D}$ 处的值等价于 $\mathrm{C}\times_{\mathrm{D}}\mathrm{D}_{d/}$。这个纤维化用于 Quillen 定理 A 的证明。
\end{example}

\begin{exercise}\label{L20-5.4}
设 $F\colon\mathrm{C}\to\mathrm{D}$ 是拟范畴之间的任意范畴纤维化。证明由 $\mathrm{D}\to\Fun(\Delta^1,\mathrm{D})$、$d\mapsto\id_d$ 诱导的态射
\[
\mathrm{C}^{\flat}\to(\Fun(\Delta^1,\mathrm{D})\times_{\mathrm{D}}\mathrm{C})^{\natural}
\]
是 $(\Set^+_\Delta)^{\mathrm{Cart}}_{/\mathrm{D}}$ 中的纤维替换。
\end{exercise}

\begin{example}\label{L20-5.5}
设 $\mathrm{C}$ 是 $\infty$-范畴。左纤维化
\[
\mathrm{TwArr}(\mathrm{C})\to\mathrm{C}^{\op}\times\mathrm{C}
\]
对应一个函子 $\mathrm{C}^{\op}\times\mathrm{C}\to\Cat_\infty$，它把 $(x,y)$ 送到一个等价于 $\mathrm{Maps}_{\mathrm{C}}(x,y)$ 的 $\infty$-群胚。
\end{example}

\begin{exercise}\label{L20-5.6}
设 $F\colon\mathrm{D}\to\mathrm{Pr}^{L}$ 是函子。复合
\[
\mathrm{D}\to\mathrm{Pr}^{L}\to\widehat{\Cat}_\infty
\]
对应一个（未必小的）余笛卡尔纤维化 $\mathrm{E}\to\mathrm{D}$。证明它同时也是笛卡尔纤维化，并且相应的函子 $\mathrm{D}^{\op}\to\widehat{\Cat}_\infty$ 穿过 $\mathrm{Pr}^{R}$ 分解。
\end{exercise}

\nitem\label{L20-5.7} 由 Grothendieck–Lurie 构造，对任意单纯集 $S$，图 $u\colon S^{\op}\to\Cat_\infty$ 如下对应到单纯集之间的笛卡尔纤维化 $X\to S$：
\begin{itemize}[leftmargin=2.2em,itemsep=0.1em]
\item 利用等价 $\Cat_\infty\simeq\mathfrak{N}((\widetilde{\Set}^+_\Delta)^\circ)$ 把 $u$ 翻译成函子
\[
\mathfrak{C}(S^{\op})\to(\widetilde{\Set}^+_\Delta)^\circ.
\]
\item 取纤维替换，并应用反拉直函子，得到 $\mathrm{Cart}_{/S}$ 中的纤维对象。
\item 设 $X\to S$ 是对应于这个纤维对象的笛卡尔纤维化。
\end{itemize}

\begin{proposition}\label{L20-5.8}（HTT.3.3.3）
设 $u\colon S^{\op}\to\Cat_\infty$ 是对应于笛卡尔纤维化 $p\colon X\to S$ 的 $\infty$-范畴图。则 $u$ 的极限可以等同于 $p$ 的笛卡尔截面的拟范畴：
\[
\lim u\simeq\operatorname{Sect}^{\mathrm{Cart}}_{/S}(X).
\]
\end{proposition}

\begin{proposition}\label{L20-5.9}（HTT.3.3.4）
设 $u\colon S\to\Cat_\infty$ 是对应于余笛卡尔纤维化 $p\colon X\to S$ 的 $\infty$-范畴图。则 $u$ 的余极限可由带标记单纯集 $X^{\natural}$ 建模，其中带标记的边是 $p$-余笛卡尔的边：
\[
\colim u\simeq X^{\natural}.
\]
\end{proposition}

\appendixsection{对偶纤维化与共轭纤维化}

\nitem\label{L20-A.1} 设 $F\colon S\to\Cat_\infty$ 是 $\infty$-范畴的图。我们可以构造下述纤维化：
\begin{enumerate}[label=(\roman*),leftmargin=2.4em,itemsep=0.1em]
\item 对应于 $F$ 的、$S^{\op}$ 上的笛卡尔纤维化；
\item 对应于 $F$ 的、$S$ 上的余笛卡尔纤维化；
\item 对应于 $\op\circ F$ 的、$S^{\op}$ 上的笛卡尔纤维化；
\item 对应于 $\op\circ F$ 的、$S$ 上的余笛卡尔纤维化。
\end{enumerate}

\begin{exercise}\label{L20-A.2}
描述如何从上述纤维化中的一个直接构造出另一个。
\end{exercise}

\nitem\label{L20-A.3} 推荐阅读。Ker.04C0。
